\newpage
\section{Conclusion and Perspectives}

\subsection{Discussion and Limitations}

\subsubsection{Project contributions}
This project has developed a comprehensive theoretical and numerical framework for analyzing opportunistic RAN sharing between mobile operators.

On the theoretical front, we demonstrated the super-additivity of the value function under the prudent rule, which guarantees the collective benefit of cooperation. We also highlighted, through numerical simulations, that this property tightly depends on the traffic distribution mechanism: it is violated under the uniform rule and satisfied under the prudent rule. This result underscores the crucial importance of choosing the right allocation rule when designing a RAN sharing system.

On the numerical front, we showed that cooperation yields an $11.9\%$ gain compared to non-cooperative operations, and that the online mode loses only $2.14\%$ of value compared to the oracle, without any capacity outages. This latter result is highly encouraging for the operational feasibility of the mechanism.

Regarding gain sharing, we proposed and compared three methods based on the Shapley value, highlighting the trade-offs between simplicity, fairness, and accounting for the specific role of guardians.

\subsubsection{Limitations}
Several limitations must be highlighted:
\begin{itemize}
    \item \textbf{Synthetic traffic:} The traffic profiles used are constructed from sinusoidal functions, which capture broad trends but do not replicate the fine-grained variability of real-world traffic (unpredictable peaks, exceptional events). Validation on real data would be necessary to confirm the robustness of the results.
    \item \textbf{Algorithmic complexity:} The approach is exponential in the number of operators. In practice, for a radio site in France ($N \le 4$), the computation is instantaneous. For extensions to multi-site or multi-cell scenarios, approximation heuristics would be required, such as Monte Carlo approximation for the Shapley value or local search algorithms for guardian optimization.
    \item \textbf{Single radio site:} The model considers an isolated radio site. In reality, operators simultaneously manage thousands of sites with interdependencies (an operator can be a guardian on one site and asleep on another). A multi-site extension would drastically increase the decision space size.
    \item \textbf{Simple prediction:} Traffic prediction relies on a simple moving average. More sophisticated methods (ARIMA, recurrent neural networks, transformers) could improve prediction quality and narrow the oracle-online gap.
    \item \textbf{Operational constraints:} The model does not account for the sleep and wake-up times of equipment, Quality of Service requirements (latency, minimum throughput), or regulatory constraints on infrastructure sharing between competing operators.
\end{itemize}

\subsubsection{Perspectives}
Several avenues for extension are possible. The use of real traffic data, provided by Orange during the internship phase, would validate the model under more realistic conditions. Furthermore, our code could be significantly improved: certain functions, such as the load function or the sharing mechanism, are currently hardcoded. It is necessary to redesign the selection mechanism for these functions to allow users to easily choose their desired strategy (e.g., via an input parameter or a dedicated interface) rather than imposing a fixed configuration within the code itself. 

Extending the model to a multi-site scenario, coupled with approximation heuristics, would pave the way for full-scale network applications. Studying dynamic sharing mechanisms, where the coalition and guardians are reconfigured in real time, would bring the model much closer to operational constraints. Finally, the formal proof of super-additivity under the uniform rule remains an open problem.

Another extension would penalize high load rates or impose explicit load
caps to account for equipment wear and operational robustness; under such
constraints, the fixed-order greedy traffic distribution would no longer
be necessarily optimal.


\subsection{Task Allocation and External Contacts}

\subsubsection{Group organization}
The work was conducted collectively, with a division of labor tailored to the different phases of the project:
\begin{itemize}
    \item \textbf{Modeling and mathematical formalism group:} Defined the theoretical framework, proved super-additivity, and designed the gain-sharing methods.
    \item \textbf{Implementation and simulation group:} Developed the Python code (allocation, optimization, and simulation modules), ran numerical tests, verified super-additivity, and generated figures.
    \item \textbf{Writing and coordination group:} Drafted the intermediate and final reports, coordinated between sub-groups, and prepared presentations.
\end{itemize}

\subsubsection{External contacts}
The project benefited from regular guidance by Azeddine Gati and Ridha Nasri, researchers at Orange, who served as tutors. Discussions focused on the model's relevance to real-world mobile network operational constraints, the choice of cost functions, and the orders of magnitude for parameters. The pedagogical coordinator, Yoan Tardy, monitored the project's progress.


\subsection{Conclusion}

This collective scientific project has successfully developed a complete framework (theoretical, algorithmic, and numerical) to analyze opportunistic RAN sharing among mobile operators.

Theoretically, we formalized the problem within the framework of coalitional game theory, introducing the concept of guardian operators and studying three traffic distribution rules. The super-additivity property, which guarantees the collective benefit of cooperation, was proven analytically under the prudent rule and verified numerically under the uniform rule. A notable negative result is that this property is violated under the uniform rule, illustrating the tight bond between the traffic allocation mechanism and the properties of the coalitional game. Nonetheless, this result does not preclude the use of this rule, as the observed violations are negligible compared to the overall gains brought by the coalition.

Numerically, simulations of a four-operator scenario over 60 time steps demonstrate a cooperative gain for the coalition compared to standalone operations. The online mode, based on simple traffic prediction with a $15\%$ safety margin, loses very little value compared to the oracle mode ($2\%$ in our simulations) without any capacity outages. This demonstrates the robustness of the mechanism against prediction uncertainties. This robustness must, of course, be contextualized, as results could change if traffic profiles evolve drastically. However, RAN sharing is primarily intended for low-traffic periods, meaning that traffic spikes during these windows would inherently remain lower.

Three gain-sharing methods were proposed and compared, highlighting the trade-offs between simplicity, fairness, and accounting for the specific role of guardians. All methods guarantee that each operator individually benefits from cooperation. While exact payoffs vary based on parameters and rules, the overall economic gain is large enough to ensure that all participating operators benefit.

Future perspectives for this work include validation on real-world traffic data, scaling up to multi-site scenarios using approximation heuristics, and investigating real-time dynamic decision mechanisms. Real-data validation remains the primary milestone, as simulation results could be challenged if real-world behavior diverges significantly from our simulated assumptions.
